% !PRINTED 8
\weberpage{8}

We too shall in what follows not renounce the aid of geometrical
intuition, and shall, for example, without hesitation regard the points
of a straight line as a continuous set.

\begin{center}
\rule{0.18\linewidth}{0.4pt}
\end{center}

An ordered set $\mathfrak M$ is called \gesperrt{measurable} under the
following assumptions: addition and multiplication are generally
executable in $\mathfrak M$, and likewise the subtraction of a smaller
element from a greater one; that is, from any two elements $a,b$ (which
may also be identical) a new element $a+b$ of $\mathfrak M$ can be
derived according to a definite rule, in such a way that $a+b$ is
greater than $a$ and greater than $b$, and the familiar laws of addition
expressed in the formulae
\[
 a+b=b+a,
 \qquad (a+b)+c=a+(b+c)
\]
hold. Moreover, for two elements $a,c$, of which the second is the
greater, a third element $b$ can be found such that $a+b=c$, which is
also expressed by the sign of subtraction $b=c-a$. Thus two unequal
elements of $\mathfrak M$ always have a definite difference. From these
assumptions it follows that a sum becomes greater when one of its
summands becomes greater. The repeated, say $m$-fold, addition of the
same element $a$ is called multiplication, and its result is denoted by
$ma$.

Finally there is added one more assumption, namely that, for every
given $a$, a sufficiently high multiple $ma$ is greater than any other
element $b$ given in advance. Among the elements of a measurable set
there is therefore no greatest element.

In a dense measurable set there is also no least element. For if $a$
were the least and $b$ an arbitrary element, then between $b$ and $b+a$
no elements could lie; for if $b<c<b+a$, then it would follow from the
definition of measurability that $a'=c-b$ would be smaller than $a$.
Conversely it also follows that a measurable set in which no least
element occurs is dense. For if $a$ and $a+b$ are any two elements, one
need only add to $a$ an element smaller than $b$ in order to obtain an
element between $a$ and $a+b$.

If a set is continuous, then the assumptions for measurability can be
simplified still further, because then subtraction is a consequence of
addition. For if $a$ and $c$ are two ele-
%CONT>>
